\documentclass[11pt,a4paper]{article}
% lesson-preamble.tex -- Shared preamble for all Phase 00 lessons
% Usage: % lesson-preamble.tex -- Shared preamble for all Phase 00 lessons
% Usage: % lesson-preamble.tex -- Shared preamble for all Phase 00 lessons
% Usage: \input{../lesson-preamble} (from subdomain directories)

% ---- Encoding and fonts ----
\usepackage[utf8]{inputenc}
\usepackage[T1]{fontenc}

% ---- Mathematics ----
\usepackage{amsmath,amssymb,amsthm}
\usepackage{mathtools}

% ---- Page geometry ----
\usepackage[margin=1in,headheight=14pt]{geometry}

% ---- Hyperlinks ----
\usepackage[colorlinks=true,linkcolor=red,citecolor=blue,urlcolor=blue]{hyperref}

% ---- Lists ----
\usepackage{enumitem}

% ---- Colors and boxes ----
\usepackage{xcolor}
\usepackage[theorems,skins,breakable]{tcolorbox}

% ---- Header/Footer ----
\usepackage{fancyhdr}
\renewcommand{\headrulewidth}{0.4pt}

% ---- Tables ----
\usepackage{booktabs}

% ---- Bibliography ----
\usepackage{natbib}
\bibliographystyle{plainnat}

% --------------------------------------------------------------------------
% Theorem environments
% --------------------------------------------------------------------------
\theoremstyle{plain}
\newtheorem{theorem}{Theorem}[section]
\newtheorem{lemma}[theorem]{Lemma}
\newtheorem{proposition}[theorem]{Proposition}
\newtheorem{corollary}[theorem]{Corollary}

\theoremstyle{definition}
\newtheorem{definition}[theorem]{Definition}
\newtheorem{example}[theorem]{Example}
\newtheorem{exercise}[theorem]{Exercise}
\newtheorem{assumption}[theorem]{Assumption}

\theoremstyle{remark}
\newtheorem{remark}[theorem]{Remark}

% --------------------------------------------------------------------------
% Colored boxes
% --------------------------------------------------------------------------
\newtcolorbox{keyresult}[1][]{%
  colback=green!5!white,
  colframe=green!50!black,
  fonttitle=\bfseries,
  title={Key Result},
  breakable,
  #1
}

\newtcolorbox{rlconnection}[1][]{%
  colback=blue!5!white,
  colframe=blue!50!black,
  fonttitle=\bfseries,
  title={RL/ML Connection},
  breakable,
  #1
}

\newtcolorbox{warning}[1][]{%
  colback=red!5!white,
  colframe=red!50!black,
  fonttitle=\bfseries,
  title={Warning},
  breakable,
  #1
}

% --------------------------------------------------------------------------
% Custom commands -- number sets
% --------------------------------------------------------------------------
\newcommand{\R}{\mathbb{R}}
\newcommand{\N}{\mathbb{N}}
\newcommand{\Q}{\mathbb{Q}}
\newcommand{\Z}{\mathbb{Z}}
\newcommand{\C}{\mathbb{C}}
\newcommand{\F}{\mathbb{F}}

% --------------------------------------------------------------------------
% Custom commands -- probability and expectation
% --------------------------------------------------------------------------
\newcommand{\E}{\mathbb{E}}
\newcommand{\Prob}{\mathbb{P}}
\newcommand{\Var}{\operatorname{Var}}
\newcommand{\Cov}{\operatorname{Cov}}
\newcommand{\indicator}{\mathbf{1}}

% --------------------------------------------------------------------------
% Custom commands -- convergence arrows
% --------------------------------------------------------------------------
\newcommand{\cas}{\xrightarrow{\text{a.s.}}}
\newcommand{\cprob}{\xrightarrow{\Prob}}
\newcommand{\clp}[1]{\xrightarrow{L^{#1}}}
\newcommand{\cdist}{\xrightarrow{d}}

% --------------------------------------------------------------------------
% Custom commands -- common operators
% --------------------------------------------------------------------------
\DeclareMathOperator{\dom}{dom}
\DeclareMathOperator{\epi}{epi}
\DeclareMathOperator{\conv}{conv}
\DeclareMathOperator{\relint}{relint}
\DeclareMathOperator{\aff}{aff}
\DeclareMathOperator{\cl}{cl}
\DeclareMathOperator{\interior}{int}
\DeclareMathOperator{\tr}{tr}
\DeclareMathOperator{\diag}{diag}
\DeclareMathOperator{\rank}{rank}
\DeclareMathOperator{\sgn}{sgn}
\DeclareMathOperator{\supp}{supp}
\DeclareMathOperator{\argmin}{arg\,min}
\DeclareMathOperator{\argmax}{arg\,max}
\DeclareMathOperator{\prox}{prox}
\DeclareMathOperator{\dist}{dist}
\DeclareMathOperator{\diam}{diam}
\DeclareMathOperator{\Span}{span}

% --------------------------------------------------------------------------
% Custom commands -- linear algebra operators
% --------------------------------------------------------------------------
\DeclareMathOperator{\nullity}{nullity}
\DeclareMathOperator{\range}{range}
\DeclareMathOperator{\nullspace}{null}
\DeclareMathOperator{\proj}{proj}

% --------------------------------------------------------------------------
% Custom commands -- measure theory
% --------------------------------------------------------------------------
\newcommand{\powerset}{\mathcal{P}}
\newcommand{\Borel}{\mathcal{B}}
 (from subdomain directories)

% ---- Encoding and fonts ----
\usepackage[utf8]{inputenc}
\usepackage[T1]{fontenc}

% ---- Mathematics ----
\usepackage{amsmath,amssymb,amsthm}
\usepackage{mathtools}

% ---- Page geometry ----
\usepackage[margin=1in,headheight=14pt]{geometry}

% ---- Hyperlinks ----
\usepackage[colorlinks=true,linkcolor=red,citecolor=blue,urlcolor=blue]{hyperref}

% ---- Lists ----
\usepackage{enumitem}

% ---- Colors and boxes ----
\usepackage{xcolor}
\usepackage[theorems,skins,breakable]{tcolorbox}

% ---- Header/Footer ----
\usepackage{fancyhdr}
\renewcommand{\headrulewidth}{0.4pt}

% ---- Tables ----
\usepackage{booktabs}

% ---- Bibliography ----
\usepackage{natbib}
\bibliographystyle{plainnat}

% --------------------------------------------------------------------------
% Theorem environments
% --------------------------------------------------------------------------
\theoremstyle{plain}
\newtheorem{theorem}{Theorem}[section]
\newtheorem{lemma}[theorem]{Lemma}
\newtheorem{proposition}[theorem]{Proposition}
\newtheorem{corollary}[theorem]{Corollary}

\theoremstyle{definition}
\newtheorem{definition}[theorem]{Definition}
\newtheorem{example}[theorem]{Example}
\newtheorem{exercise}[theorem]{Exercise}
\newtheorem{assumption}[theorem]{Assumption}

\theoremstyle{remark}
\newtheorem{remark}[theorem]{Remark}

% --------------------------------------------------------------------------
% Colored boxes
% --------------------------------------------------------------------------
\newtcolorbox{keyresult}[1][]{%
  colback=green!5!white,
  colframe=green!50!black,
  fonttitle=\bfseries,
  title={Key Result},
  breakable,
  #1
}

\newtcolorbox{rlconnection}[1][]{%
  colback=blue!5!white,
  colframe=blue!50!black,
  fonttitle=\bfseries,
  title={RL/ML Connection},
  breakable,
  #1
}

\newtcolorbox{warning}[1][]{%
  colback=red!5!white,
  colframe=red!50!black,
  fonttitle=\bfseries,
  title={Warning},
  breakable,
  #1
}

% --------------------------------------------------------------------------
% Custom commands -- number sets
% --------------------------------------------------------------------------
\newcommand{\R}{\mathbb{R}}
\newcommand{\N}{\mathbb{N}}
\newcommand{\Q}{\mathbb{Q}}
\newcommand{\Z}{\mathbb{Z}}
\newcommand{\C}{\mathbb{C}}
\newcommand{\F}{\mathbb{F}}

% --------------------------------------------------------------------------
% Custom commands -- probability and expectation
% --------------------------------------------------------------------------
\newcommand{\E}{\mathbb{E}}
\newcommand{\Prob}{\mathbb{P}}
\newcommand{\Var}{\operatorname{Var}}
\newcommand{\Cov}{\operatorname{Cov}}
\newcommand{\indicator}{\mathbf{1}}

% --------------------------------------------------------------------------
% Custom commands -- convergence arrows
% --------------------------------------------------------------------------
\newcommand{\cas}{\xrightarrow{\text{a.s.}}}
\newcommand{\cprob}{\xrightarrow{\Prob}}
\newcommand{\clp}[1]{\xrightarrow{L^{#1}}}
\newcommand{\cdist}{\xrightarrow{d}}

% --------------------------------------------------------------------------
% Custom commands -- common operators
% --------------------------------------------------------------------------
\DeclareMathOperator{\dom}{dom}
\DeclareMathOperator{\epi}{epi}
\DeclareMathOperator{\conv}{conv}
\DeclareMathOperator{\relint}{relint}
\DeclareMathOperator{\aff}{aff}
\DeclareMathOperator{\cl}{cl}
\DeclareMathOperator{\interior}{int}
\DeclareMathOperator{\tr}{tr}
\DeclareMathOperator{\diag}{diag}
\DeclareMathOperator{\rank}{rank}
\DeclareMathOperator{\sgn}{sgn}
\DeclareMathOperator{\supp}{supp}
\DeclareMathOperator{\argmin}{arg\,min}
\DeclareMathOperator{\argmax}{arg\,max}
\DeclareMathOperator{\prox}{prox}
\DeclareMathOperator{\dist}{dist}
\DeclareMathOperator{\diam}{diam}
\DeclareMathOperator{\Span}{span}

% --------------------------------------------------------------------------
% Custom commands -- linear algebra operators
% --------------------------------------------------------------------------
\DeclareMathOperator{\nullity}{nullity}
\DeclareMathOperator{\range}{range}
\DeclareMathOperator{\nullspace}{null}
\DeclareMathOperator{\proj}{proj}

% --------------------------------------------------------------------------
% Custom commands -- measure theory
% --------------------------------------------------------------------------
\newcommand{\powerset}{\mathcal{P}}
\newcommand{\Borel}{\mathcal{B}}
 (from subdomain directories)

% ---- Encoding and fonts ----
\usepackage[utf8]{inputenc}
\usepackage[T1]{fontenc}

% ---- Mathematics ----
\usepackage{amsmath,amssymb,amsthm}
\usepackage{mathtools}

% ---- Page geometry ----
\usepackage[margin=1in,headheight=14pt]{geometry}

% ---- Hyperlinks ----
\usepackage[colorlinks=true,linkcolor=red,citecolor=blue,urlcolor=blue]{hyperref}

% ---- Lists ----
\usepackage{enumitem}

% ---- Colors and boxes ----
\usepackage{xcolor}
\usepackage[theorems,skins,breakable]{tcolorbox}

% ---- Header/Footer ----
\usepackage{fancyhdr}
\renewcommand{\headrulewidth}{0.4pt}

% ---- Tables ----
\usepackage{booktabs}

% ---- Bibliography ----
\usepackage{natbib}
\bibliographystyle{plainnat}

% --------------------------------------------------------------------------
% Theorem environments
% --------------------------------------------------------------------------
\theoremstyle{plain}
\newtheorem{theorem}{Theorem}[section]
\newtheorem{lemma}[theorem]{Lemma}
\newtheorem{proposition}[theorem]{Proposition}
\newtheorem{corollary}[theorem]{Corollary}

\theoremstyle{definition}
\newtheorem{definition}[theorem]{Definition}
\newtheorem{example}[theorem]{Example}
\newtheorem{exercise}[theorem]{Exercise}
\newtheorem{assumption}[theorem]{Assumption}

\theoremstyle{remark}
\newtheorem{remark}[theorem]{Remark}

% --------------------------------------------------------------------------
% Colored boxes
% --------------------------------------------------------------------------
\newtcolorbox{keyresult}[1][]{%
  colback=green!5!white,
  colframe=green!50!black,
  fonttitle=\bfseries,
  title={Key Result},
  breakable,
  #1
}

\newtcolorbox{rlconnection}[1][]{%
  colback=blue!5!white,
  colframe=blue!50!black,
  fonttitle=\bfseries,
  title={RL/ML Connection},
  breakable,
  #1
}

\newtcolorbox{warning}[1][]{%
  colback=red!5!white,
  colframe=red!50!black,
  fonttitle=\bfseries,
  title={Warning},
  breakable,
  #1
}

% --------------------------------------------------------------------------
% Custom commands -- number sets
% --------------------------------------------------------------------------
\newcommand{\R}{\mathbb{R}}
\newcommand{\N}{\mathbb{N}}
\newcommand{\Q}{\mathbb{Q}}
\newcommand{\Z}{\mathbb{Z}}
\newcommand{\C}{\mathbb{C}}
\newcommand{\F}{\mathbb{F}}

% --------------------------------------------------------------------------
% Custom commands -- probability and expectation
% --------------------------------------------------------------------------
\newcommand{\E}{\mathbb{E}}
\newcommand{\Prob}{\mathbb{P}}
\newcommand{\Var}{\operatorname{Var}}
\newcommand{\Cov}{\operatorname{Cov}}
\newcommand{\indicator}{\mathbf{1}}

% --------------------------------------------------------------------------
% Custom commands -- convergence arrows
% --------------------------------------------------------------------------
\newcommand{\cas}{\xrightarrow{\text{a.s.}}}
\newcommand{\cprob}{\xrightarrow{\Prob}}
\newcommand{\clp}[1]{\xrightarrow{L^{#1}}}
\newcommand{\cdist}{\xrightarrow{d}}

% --------------------------------------------------------------------------
% Custom commands -- common operators
% --------------------------------------------------------------------------
\DeclareMathOperator{\dom}{dom}
\DeclareMathOperator{\epi}{epi}
\DeclareMathOperator{\conv}{conv}
\DeclareMathOperator{\relint}{relint}
\DeclareMathOperator{\aff}{aff}
\DeclareMathOperator{\cl}{cl}
\DeclareMathOperator{\interior}{int}
\DeclareMathOperator{\tr}{tr}
\DeclareMathOperator{\diag}{diag}
\DeclareMathOperator{\rank}{rank}
\DeclareMathOperator{\sgn}{sgn}
\DeclareMathOperator{\supp}{supp}
\DeclareMathOperator{\argmin}{arg\,min}
\DeclareMathOperator{\argmax}{arg\,max}
\DeclareMathOperator{\prox}{prox}
\DeclareMathOperator{\dist}{dist}
\DeclareMathOperator{\diam}{diam}
\DeclareMathOperator{\Span}{span}

% --------------------------------------------------------------------------
% Custom commands -- linear algebra operators
% --------------------------------------------------------------------------
\DeclareMathOperator{\nullity}{nullity}
\DeclareMathOperator{\range}{range}
\DeclareMathOperator{\nullspace}{null}
\DeclareMathOperator{\proj}{proj}

% --------------------------------------------------------------------------
% Custom commands -- measure theory
% --------------------------------------------------------------------------
\newcommand{\powerset}{\mathcal{P}}
\newcommand{\Borel}{\mathcal{B}}

\pagestyle{fancy}
\fancyhf{}
\fancyhead[L]{\textsc{Lesson 12: Inequalities and Limit Theorems}}
\fancyhead[R]{\thepage}
\title{Lesson 12: Inequalities and Limit Theorems}
\author{AI/ML Mathematics}
\date{}
\begin{document}
\maketitle
\tableofcontents
\newpage

%% ============================================================
%% SECTION 1: MOTIVATION AND OVERVIEW
%% ============================================================

\section{Motivation and Overview}
\label{sec:motivation}

Suppose a reinforcement learning agent uses Monte Carlo rollouts to estimate
the value of a policy: it runs $n$ independent episodes starting from a
fixed state $s$, collects the discounted returns
$G^{(1)}, \ldots, G^{(n)}$, and forms the sample mean
$\bar{G}_n = \frac{1}{n}\sum_{i=1}^n G^{(i)}$. The true value
$V^\pi(s) = \E[G^{(i)}]$ is unknown. Three fundamental questions arise:
\begin{enumerate}[nosep]
    \item Does $\bar{G}_n$ converge to $V^\pi(s)$ as $n \to \infty$,
    and in what sense?
    \item How quickly does $\bar{G}_n$ concentrate around $V^\pi(s)$ --
    that is, how large must $n$ be so that
    $|\bar{G}_n - V^\pi(s)| < \varepsilon$ with probability at least
    $1 - \delta$?
    \item What is the approximate distribution of the estimation error
    $\bar{G}_n - V^\pi(s)$ for large $n$?
\end{enumerate}
The first question is answered by the law of large numbers, the second by
concentration inequalities (Chebyshev and Hoeffding), and the third by the
central limit theorem. The tools developed in this lesson will let us
prove these results, derive explicit sample complexity bounds for Monte
Carlo estimation, and understand why the normal distribution appears as
a universal limit.

\begin{rlconnection}[title={Why Inequalities and Limit Theorems
    Matter for RL/ML}]
\begin{itemize}[nosep]
    \item \textbf{Monte Carlo policy evaluation}: The law of large
    numbers guarantees that averaging returns converges to the true
    value function; concentration inequalities quantify how many
    episodes are needed for a given accuracy.
    \item \textbf{TD-learning convergence rates}: Stochastic
    approximation convergence proofs rely on Markov's and Chebyshev's
    inequalities to control the error at each iteration (Phase~01,
    Lesson on Stochastic Approximation).
    \item \textbf{PAC-MDP bounds}: Sample complexity results in
    PAC-MDP theory use Hoeffding's inequality to bound the
    probability of inaccurate value estimates.
    \item \textbf{Bandit confidence intervals}: UCB algorithms
    construct confidence bounds for arm means using Hoeffding's
    inequality; the central limit theorem justifies Gaussian
    approximations for large sample sizes.
\end{itemize}
\end{rlconnection}

\paragraph{Prerequisites.}
Probability spaces and the Kolmogorov axioms (Lesson~9). Random
variables, PMFs, PDFs, CDFs, and standard distributions -- in
particular the Bernoulli, Poisson, and Normal distributions
(Lesson~10). Expectation, variance, covariance, independence of
random variables, the law of the unconscious statistician (LOTUS),
linearity of expectation, and the tower property (Lesson~11).

\paragraph{Learning objectives.}
After completing this lesson, the student will be able to:
\begin{itemize}
    \item \emph{Define} the moment generating function and \emph{derive}
    moments from it by differentiation.
    \item \emph{Prove} Markov's, Chebyshev's, Cauchy-Schwarz, and
    Jensen's inequalities.
    \item \emph{State} the definitions of convergence in probability
    and convergence in distribution, and \emph{prove} the weak law of
    large numbers via Chebyshev's inequality.
    \item \emph{State} the strong law of large numbers and the central
    limit theorem, and \emph{apply} them in worked examples.
    \item \emph{Prove} Hoeffding's lemma and Hoeffding's inequality,
    and \emph{derive} sample complexity bounds for Monte Carlo
    estimation.
    \item \emph{Compare} the sample complexity guarantees provided by
    Chebyshev's inequality versus Hoeffding's inequality.
\end{itemize}

%% ============================================================
%% SECTION 2: REFERENCES
%% ============================================================

\section{References}
\label{sec:references}

\begin{itemize}
    \item \citet[Chapters 7--8]{ross2019} -- Moment generating
    functions, limit theorems, and concentration inequalities.
    \item \citet[Chapter 10]{blitzstein2019} -- Inequalities and
    limit theorems with extensive worked examples and RL-relevant
    applications.
    \item \citet[Chapters 8--10]{feller1968} -- Classical treatment
    of the law of large numbers and the central limit theorem.
    \item \citet[Sections 2.2--2.3]{durrett2019} -- Rigorous
    treatment of the weak and strong laws of large numbers and the
    central limit theorem via characteristic functions.
    \item \citet[Sections 5.1--5.4]{grimmett2020} -- Generating
    functions, moment generating functions, and convergence of random
    variables.
\end{itemize}

%% ============================================================
%% SECTION 3: MOMENT GENERATING FUNCTIONS
%% ============================================================

\section{Moment Generating Functions}
\label{sec:mgf}

Moment generating functions (MGFs) provide a powerful technique for
computing moments, identifying distributions, and establishing the
distribution of sums of independent random variables.

\begin{definition}[Moment Generating Function]
\label{def:mgf}
Let $X$ be a random variable. The \emph{moment generating function}
(MGF) of $X$ is the function $M_X \colon \R \to [0, \infty]$ defined by
\[
    M_X(t) = \E[e^{tX}],
\]
provided the expectation exists. We say the MGF exists in a
neighborhood of the origin if there exists $h > 0$ such that
$M_X(t) < \infty$ for all $|t| < h$.
\end{definition}

\begin{remark}
\label{rem:mgf_exists}
Note that $M_X(0) = \E[e^0] = 1$ always holds. The MGF may fail to
exist for $t \neq 0$; for instance, the Cauchy distribution has no
finite moments beyond order zero, so its MGF is infinite for all
$t \neq 0$. When we say ``the MGF exists,'' we mean it is finite in
some open interval containing the origin.
\end{remark}

\begin{theorem}[Moments from the MGF]
\label{thm:moments_from_mgf}
If the MGF $M_X(t)$ exists in a neighborhood of the origin, then
$\E[|X|^n] < \infty$ for all $n \geq 1$ and
\[
    \E[X^n] = M_X^{(n)}(0),
\]
where $M_X^{(n)}$ denotes the $n$-th derivative of $M_X$.
\end{theorem}

\begin{proof}
Since $M_X(t) = \E[e^{tX}]$ is finite for $|t| < h$, we can
differentiate under the expectation sign (justified by the dominated
convergence theorem, since $|X^n e^{tX}|$ is dominated by a sum of
exponentials that are integrable for $|t|$ sufficiently small). We
have
\[
    M_X'(t) = \frac{d}{dt} \E[e^{tX}] = \E\!\left[\frac{d}{dt} e^{tX}\right]
    = \E[X e^{tX}].
\]
Evaluating at $t = 0$: $M_X'(0) = \E[X \cdot 1] = \E[X]$.

More generally, by induction, the $n$-th derivative is
\[
    M_X^{(n)}(t) = \E[X^n e^{tX}],
\]
and evaluating at $t = 0$ gives $M_X^{(n)}(0) = \E[X^n]$.

For the base case $n = 1$, we have shown $M_X'(t) = \E[Xe^{tX}]$.
Assume $M_X^{(k)}(t) = \E[X^k e^{tX}]$ for some $k \geq 1$. Then
\[
    M_X^{(k+1)}(t) = \frac{d}{dt} \E[X^k e^{tX}]
    = \E\!\left[X^k \frac{d}{dt} e^{tX}\right]
    = \E[X^{k+1} e^{tX}],
\]
completing the induction.
\end{proof}

\begin{example}[First Two Moments via the MGF]
\label{ex:moments_via_mgf}
The first two moments are $\E[X] = M_X'(0)$ and
$\E[X^2] = M_X''(0)$. The variance can then be computed as
\[
    \Var(X) = \E[X^2] - (\E[X])^2 = M_X''(0) - (M_X'(0))^2.
\]
\end{example}

\subsection{MGF of a Sum of Independent Random Variables}

\begin{theorem}[MGF of a Sum of Independents]
\label{thm:mgf_sum}
If $X$ and $Y$ are independent random variables whose MGFs exist,
then
\[
    M_{X+Y}(t) = M_X(t) \cdot M_Y(t).
\]
More generally, if $X_1, \ldots, X_n$ are independent, then
$M_{X_1 + \cdots + X_n}(t) = \prod_{i=1}^n M_{X_i}(t)$.
\end{theorem}

\begin{proof}
By definition,
\[
    M_{X+Y}(t) = \E[e^{t(X+Y)}] = \E[e^{tX} e^{tY}].
\]
Since $X$ and $Y$ are independent, $e^{tX}$ and $e^{tY}$ are also
independent (as they are functions of independent random variables;
see Lesson~11, Proposition~\ref*{prop:indep_uncorrelated}). Therefore
\[
    \E[e^{tX} e^{tY}] = \E[e^{tX}] \cdot \E[e^{tY}]
    = M_X(t) \cdot M_Y(t).
\]
The general case follows by induction on $n$.
\end{proof}

\begin{remark}
\label{rem:mgf_uniqueness}
A fundamental result from the theory of Laplace transforms (see,
e.g., \citet[Theorem 2.3.6]{grimmett2020}) states that if the MGF
of $X$ exists in a neighborhood of the origin, then it uniquely
determines the distribution of $X$. That is, if $M_X(t) = M_Y(t)$
for all $t$ in some interval $(-h, h)$ with $h > 0$, then $X$ and
$Y$ have the same distribution. We state this without proof, as it
requires complex analysis (the inversion formula for Laplace
transforms).
\end{remark}

\subsection{MGFs of Standard Distributions}

\begin{proposition}[MGF of the Bernoulli Distribution]
\label{prop:mgf_bernoulli}
If $X \sim \mathrm{Bernoulli}(p)$, then
\[
    M_X(t) = 1 - p + pe^t \quad \text{for all } t \in \R.
\]
\end{proposition}

\begin{proof}
Since $X$ takes value $1$ with probability $p$ and $0$ with
probability $1-p$,
\[
    M_X(t) = \E[e^{tX}] = e^{t \cdot 0}(1-p) + e^{t \cdot 1} p
    = (1-p) + pe^t. \qedhere
\]
\end{proof}

\begin{proposition}[MGF of the Poisson Distribution]
\label{prop:mgf_poisson}
If $X \sim \mathrm{Poisson}(\lambda)$, then
\[
    M_X(t) = e^{\lambda(e^t - 1)} \quad \text{for all } t \in \R.
\]
\end{proposition}

\begin{proof}
Using the PMF $\Prob(X = k) = e^{-\lambda} \lambda^k / k!$,
\begin{align*}
    M_X(t) &= \E[e^{tX}] = \sum_{k=0}^\infty e^{tk}
    \frac{e^{-\lambda} \lambda^k}{k!}
    = e^{-\lambda} \sum_{k=0}^\infty \frac{(\lambda e^t)^k}{k!}
    = e^{-\lambda} \cdot e^{\lambda e^t}
    = e^{\lambda(e^t - 1)},
\end{align*}
where we recognized the Taylor series
$\sum_{k=0}^\infty z^k / k! = e^z$ with $z = \lambda e^t$.
\end{proof}

\begin{proposition}[MGF of the Normal Distribution]
\label{prop:mgf_normal}
If $X \sim \mathcal{N}(\mu, \sigma^2)$, then
\[
    M_X(t) = e^{\mu t + \sigma^2 t^2 / 2}
    \quad \text{for all } t \in \R.
\]
\end{proposition}

\begin{proof}
First consider $Z \sim \mathcal{N}(0, 1)$. Then
\begin{align*}
    M_Z(t) &= \E[e^{tZ}]
    = \int_{-\infty}^{\infty} e^{tz} \frac{1}{\sqrt{2\pi}}
    e^{-z^2/2} \, dz
    = \frac{1}{\sqrt{2\pi}} \int_{-\infty}^{\infty}
    e^{tz - z^2/2} \, dz.
\end{align*}
Complete the square in the exponent:
$tz - z^2/2 = -(z^2 - 2tz)/2 = -((z-t)^2 - t^2)/2
= -(z-t)^2/2 + t^2/2$.
Therefore
\begin{align*}
    M_Z(t) &= \frac{e^{t^2/2}}{\sqrt{2\pi}}
    \int_{-\infty}^{\infty} e^{-(z-t)^2/2} \, dz
    = e^{t^2/2},
\end{align*}
since $\int_{-\infty}^{\infty} e^{-(z-t)^2/2}\,dz = \sqrt{2\pi}$
by the Gaussian integral (Lesson~10) with a shift of variable
$u = z - t$.

For general $X = \mu + \sigma Z$ where $Z \sim \mathcal{N}(0,1)$,
\[
    M_X(t) = \E[e^{t(\mu + \sigma Z)}]
    = e^{\mu t} \E[e^{(t\sigma) Z}]
    = e^{\mu t} M_Z(t\sigma)
    = e^{\mu t} e^{\sigma^2 t^2/2}
    = e^{\mu t + \sigma^2 t^2/2}. \qedhere
\]
\end{proof}

\subsection{Applications of the MGF}

\begin{example}[Sum of Independent Poissons]
\label{ex:sum_poisson}
Let $X_1 \sim \mathrm{Poisson}(\lambda_1)$ and
$X_2 \sim \mathrm{Poisson}(\lambda_2)$ be independent. By
Theorem~\ref{thm:mgf_sum},
\[
    M_{X_1 + X_2}(t) = M_{X_1}(t) M_{X_2}(t)
    = e^{\lambda_1(e^t - 1)} e^{\lambda_2(e^t - 1)}
    = e^{(\lambda_1 + \lambda_2)(e^t - 1)}.
\]
This is the MGF of $\mathrm{Poisson}(\lambda_1 + \lambda_2)$.
By the uniqueness of the MGF (Remark~\ref{rem:mgf_uniqueness}),
$X_1 + X_2 \sim \mathrm{Poisson}(\lambda_1 + \lambda_2)$.
\end{example}

\begin{example}[Sum of Independent Normals]
\label{ex:sum_normal}
Let $X_1 \sim \mathcal{N}(\mu_1, \sigma_1^2)$ and
$X_2 \sim \mathcal{N}(\mu_2, \sigma_2^2)$ be independent. Then
\begin{align*}
    M_{X_1 + X_2}(t) &= e^{\mu_1 t + \sigma_1^2 t^2/2}
    \cdot e^{\mu_2 t + \sigma_2^2 t^2/2}
    = e^{(\mu_1 + \mu_2)t + (\sigma_1^2 + \sigma_2^2)t^2/2}.
\end{align*}
This is the MGF of
$\mathcal{N}(\mu_1 + \mu_2, \sigma_1^2 + \sigma_2^2)$. By
uniqueness, $X_1 + X_2 \sim
\mathcal{N}(\mu_1 + \mu_2, \sigma_1^2 + \sigma_2^2)$.
\end{example}

\begin{rlconnection}[title={MGFs and Reward Distributions}]
In RL, rewards are often modeled as sums of independent random
variables. The MGF factorization theorem
(Theorem~\ref{thm:mgf_sum}) immediately identifies the distribution
of cumulative rewards when individual rewards are independent. For
instance, if each of $n$ independent episodes yields a Poisson-distributed
count of some event with parameter $\lambda$, the total count across
all episodes is $\mathrm{Poisson}(n\lambda)$. Similarly, if reward
noise is Gaussian, the average reward over $n$ episodes remains
Gaussian with variance shrinking as $1/n$.
\end{rlconnection}

We now turn to probability inequalities, which provide deterministic
upper bounds on tail probabilities without requiring knowledge of the
full distribution.

%% ============================================================
%% SECTION 4: PROBABILITY INEQUALITIES
%% ============================================================

\section{Probability Inequalities}
\label{sec:inequalities}

Probability inequalities bound the probability that a random variable
deviates from its mean. These bounds are essential for establishing
convergence rates and deriving sample complexity results.

\subsection{Markov's Inequality}

\begin{theorem}[Markov's Inequality]
\label{thm:markov}
Let $X$ be a non-negative random variable with finite expectation.
Then for any $a > 0$,
\[
    \Prob(X \geq a) \leq \frac{\E[X]}{a}.
\]
\end{theorem}

\begin{proof}
Define the indicator $\indicator_{\{X \geq a\}}$. Since $X \geq 0$,
we have the pointwise inequality
\[
    X \geq a \cdot \indicator_{\{X \geq a\}},
\]
because when $X \geq a$ the right side equals $a \leq X$, and when
$X < a$ the right side equals $0 \leq X$. Taking expectations and
using linearity:
\[
    \E[X] \geq a \cdot \E[\indicator_{\{X \geq a\}}]
    = a \cdot \Prob(X \geq a).
\]
Dividing both sides by $a > 0$ gives
$\Prob(X \geq a) \leq \E[X]/a$.
\end{proof}

\begin{example}[Markov Applied to a Uniform Variable]
\label{ex:markov_uniform}
Let $X \sim \mathrm{Uniform}(0, 1)$, so $\E[X] = 1/2$. For
$a = 3/4$, Markov gives
$\Prob(X \geq 3/4) \leq (1/2)/(3/4) = 2/3$.
The exact probability is $\Prob(X \geq 3/4) = 1/4$, so the bound
is loose but valid. Markov's inequality uses only the mean, so its
generality comes at the cost of tightness.
\end{example}

\subsection{Chebyshev's Inequality}

\begin{theorem}[Chebyshev's Inequality]
\label{thm:chebyshev}
Let $X$ be a random variable with finite mean $\mu = \E[X]$ and
finite variance $\sigma^2 = \Var(X)$. Then for any $k > 0$,
\[
    \Prob(|X - \mu| \geq k) \leq \frac{\sigma^2}{k^2}.
\]
\end{theorem}

\begin{proof}
Apply Markov's inequality (Theorem~\ref{thm:markov}) to the
non-negative random variable $Y = (X - \mu)^2$ with threshold
$a = k^2$:
\[
    \Prob(|X - \mu| \geq k)
    = \Prob((X - \mu)^2 \geq k^2)
    \leq \frac{\E[(X - \mu)^2]}{k^2}
    = \frac{\Var(X)}{k^2}
    = \frac{\sigma^2}{k^2}. \qedhere
\]
\end{proof}

\begin{example}[Chebyshev for a Poisson Variable]
\label{ex:chebyshev_poisson}
Let $X \sim \mathrm{Poisson}(\lambda)$ with $\E[X] = \lambda$ and
$\Var(X) = \lambda$ (Lesson~10). Chebyshev gives
\[
    \Prob(|X - \lambda| \geq k) \leq \frac{\lambda}{k^2}.
\]
For $\lambda = 100$ and $k = 20$:
$\Prob(|X - 100| \geq 20) \leq 100/400 = 1/4$.
\end{example}

\begin{example}[Chebyshev for a Sample Mean]
\label{ex:chebyshev_mean}
Let $X_1, \ldots, X_n$ be i.i.d.\ with mean $\mu$ and variance
$\sigma^2$. The sample mean $\bar{X}_n = \frac{1}{n}\sum_{i=1}^n X_i$
has $\E[\bar{X}_n] = \mu$ and
$\Var(\bar{X}_n) = \sigma^2/n$ (by independence and the variance
of a sum; see Lesson~11, Corollary~4.9).
Chebyshev gives
\[
    \Prob(|\bar{X}_n - \mu| \geq \varepsilon)
    \leq \frac{\sigma^2}{n \varepsilon^2}.
\]
This bound decreases as $O(1/n)$, confirming that sample means
concentrate around the population mean.
\end{example}

\subsection{Cauchy-Schwarz Inequality for Expectations}

\begin{theorem}[Cauchy-Schwarz Inequality]
\label{thm:cauchy_schwarz}
For any random variables $X$ and $Y$ with finite second moments,
\[
    |\E[XY]|^2 \leq \E[X^2] \, \E[Y^2].
\]
Equality holds if and only if $Y = cX$ almost surely for some
constant $c \in \R$ (or $X = 0$ a.s.).
\end{theorem}

\begin{proof}
If $\E[X^2] = 0$, then $X = 0$ a.s., so $\E[XY] = 0$ and the
inequality holds trivially. Assume $\E[X^2] > 0$.

For any $t \in \R$, the random variable $(tX + Y)^2$ is non-negative,
so
\[
    0 \leq \E[(tX + Y)^2] = t^2 \E[X^2] + 2t \E[XY] + \E[Y^2].
\]
This is a quadratic in $t$ that is non-negative for all $t$. A
real quadratic $at^2 + bt + c \geq 0$ for all $t$ if and only if
its discriminant satisfies $b^2 - 4ac \leq 0$. Here
$a = \E[X^2]$, $b = 2\E[XY]$, $c = \E[Y^2]$, so
\[
    4(\E[XY])^2 - 4\E[X^2]\E[Y^2] \leq 0,
\]
which gives $(\E[XY])^2 \leq \E[X^2]\E[Y^2]$.

Equality holds if and only if the discriminant is zero, meaning
the quadratic has a real root $t_0$ with $\E[(t_0 X + Y)^2] = 0$,
so $Y = -t_0 X$ almost surely.
\end{proof}

\begin{remark}
\label{rem:cs_correlation}
The Cauchy-Schwarz inequality immediately implies the correlation
bound $|\rho(X,Y)| \leq 1$ (Lesson~11, Proposition~4.14). It also
appears in Phase~01 (Lesson~7) as a fundamental property of inner
product spaces, where $\E[XY]$ plays the role of the inner product
on $L^2$.
\end{remark}

\subsection{Jensen's Inequality}

\begin{definition}[Convex Function]
\label{def:convex}
A function $g \colon \R \to \R$ is \emph{convex} if for all
$x, y \in \R$ and $\lambda \in [0, 1]$,
\[
    g(\lambda x + (1 - \lambda) y) \leq \lambda g(x)
    + (1 - \lambda) g(y).
\]
\end{definition}

\begin{theorem}[Jensen's Inequality]
\label{thm:jensen}
Let $X$ be a random variable with finite expectation and
$g \colon \R \to \R$ a convex function. Then
\[
    g(\E[X]) \leq \E[g(X)],
\]
provided $\E[|g(X)|] < \infty$.
\end{theorem}

\begin{proof}
We prove the result for a discrete random variable $X$ taking
finitely many values, using induction on the number of values.

\textbf{Base case} ($n = 1$): If $X = x_1$ with probability $1$,
then $g(\E[X]) = g(x_1) = \E[g(X)]$, so equality holds.

\textbf{Base case} ($n = 2$): If $X$ takes values $x_1, x_2$ with
probabilities $p_1, p_2$ where $p_1 + p_2 = 1$, then
$\E[X] = p_1 x_1 + p_2 x_2$ and
\[
    g(\E[X]) = g(p_1 x_1 + p_2 x_2)
    \leq p_1 g(x_1) + p_2 g(x_2) = \E[g(X)],
\]
by the definition of convexity with $\lambda = p_1$.

\textbf{Inductive step}: Assume the result holds for random variables
taking at most $n - 1$ values. Suppose $X$ takes values
$x_1, \ldots, x_n$ with probabilities $p_1, \ldots, p_n$, where
$p_n < 1$ (otherwise $X$ is constant and the result is trivial).
Let $q = 1 - p_n > 0$. Then
\[
    \E[X] = \sum_{i=1}^n p_i x_i
    = q \underbrace{\left(\sum_{i=1}^{n-1} \frac{p_i}{q} x_i\right)}_{=: \,\mu}
    + p_n x_n.
\]
Since $\sum_{i=1}^{n-1} p_i/q = 1$, the quantity $\mu$ is the
expectation of a random variable $Y$ taking values
$x_1, \ldots, x_{n-1}$ with probabilities $p_1/q, \ldots, p_{n-1}/q$.
By convexity of $g$ (with $\lambda = q$):
\[
    g(\E[X]) = g(q\mu + p_n x_n)
    \leq q \, g(\mu) + p_n \, g(x_n).
\]
By the inductive hypothesis applied to $Y$:
$g(\mu) = g(\E[Y]) \leq \E[g(Y)]
= \sum_{i=1}^{n-1} (p_i/q) g(x_i)$.
Therefore
\[
    g(\E[X]) \leq q \sum_{i=1}^{n-1} \frac{p_i}{q} g(x_i)
    + p_n g(x_n)
    = \sum_{i=1}^n p_i g(x_i) = \E[g(X)]. \qedhere
\]
\end{proof}

\begin{example}[Jensen Applied to Common Functions]
\label{ex:jensen_examples}
\begin{enumerate}[nosep,label=(\alph*)]
    \item Since $g(x) = x^2$ is convex, Jensen gives
    $(\E[X])^2 \leq \E[X^2]$, which is equivalent to $\Var(X) \geq 0$.
    \item Since $g(x) = e^x$ is convex, Jensen gives
    $e^{\E[X]} \leq \E[e^X] = M_X(1)$.
    \item Since $g(x) = -\ln x$ is convex on $(0, \infty)$,
    Jensen gives $-\ln(\E[X]) \leq \E[-\ln X]$, i.e.,
    $\ln(\E[X]) \geq \E[\ln X]$. This is the AM-GM inequality
    in its probabilistic form.
\end{enumerate}
\end{example}

\begin{keyresult}[title={Probability Inequalities Summary}]
\begin{itemize}[nosep]
    \item \textbf{Markov}: $\Prob(X \geq a) \leq \E[X]/a$ for
    $X \geq 0$, $a > 0$. Uses only the first moment.
    \item \textbf{Chebyshev}: $\Prob(|X - \mu| \geq k) \leq
    \sigma^2/k^2$. Uses first two moments.
    \item \textbf{Cauchy-Schwarz}: $|\E[XY]|^2 \leq \E[X^2]\E[Y^2]$.
    \item \textbf{Jensen}: $g(\E[X]) \leq \E[g(X)]$ for convex $g$.
\end{itemize}
These inequalities form a hierarchy: Chebyshev follows from Markov,
and each successive result incorporates more information about the
distribution, yielding tighter bounds.
\end{keyresult}

\begin{rlconnection}[title={Concentration Inequalities and
    Convergence Rates in RL}]
Concentration inequalities translate directly into convergence rate
guarantees for RL algorithms. Chebyshev's inequality applied to the
sample mean gives an $O(1/n)$ bound on the probability of
$\varepsilon$-deviation, leading to $O(1/(\varepsilon^2 \delta))$
sample complexity (Section~\ref{sec:tail_bounds}). Exponential
inequalities such as Hoeffding's (Section~\ref{sec:tail_bounds})
improve this to $O(\ln(1/\delta)/\varepsilon^2)$, which is the
standard dependence in PAC-MDP theory and bandit algorithms.
\end{rlconnection}

With the fundamental inequalities in hand, we now study the asymptotic
behavior of sample means through the law of large numbers and the
central limit theorem.

%% ============================================================
%% SECTION 5: CONVERGENCE CONCEPTS AND THE LAW OF LARGE NUMBERS
%% ============================================================

\section{Convergence Concepts and the Law of Large Numbers}
\label{sec:lln}

\subsection{Modes of Convergence}

\begin{definition}[Convergence in Probability]
\label{def:conv_prob}
A sequence of random variables $(X_n)_{n \geq 1}$ \emph{converges in
probability} to a random variable $X$, written $X_n \cprob X$, if
for every $\varepsilon > 0$,
\[
    \lim_{n \to \infty} \Prob(|X_n - X| \geq \varepsilon) = 0.
\]
\end{definition}

\begin{definition}[Almost Sure Convergence]
\label{def:conv_as}
A sequence $(X_n)_{n \geq 1}$ \emph{converges almost surely} to $X$,
written $X_n \cas X$, if
\[
    \Prob\!\left(\lim_{n \to \infty} X_n = X\right) = 1.
\]
\end{definition}

\begin{definition}[Convergence in Distribution]
\label{def:conv_dist}
A sequence $(X_n)_{n \geq 1}$ \emph{converges in distribution} to $X$,
written $X_n \cdist X$, if
\[
    \lim_{n \to \infty} F_{X_n}(x) = F_X(x)
\]
at every point $x$ where $F_X$ is continuous. Here $F_{X_n}$ and
$F_X$ are the CDFs of $X_n$ and $X$, respectively.
\end{definition}

\begin{remark}
\label{rem:convergence_hierarchy}
Almost sure convergence implies convergence in probability, which in
turn implies convergence in distribution:
\[
    X_n \cas X \;\Longrightarrow\; X_n \cprob X
    \;\Longrightarrow\; X_n \cdist X.
\]
The converses are generally false. We will not prove these
implications here; they are established rigorously in Phase~01
(measure-theoretic probability).
\end{remark}

\subsection{Weak Law of Large Numbers}

\begin{theorem}[Weak Law of Large Numbers (WLLN)]
\label{thm:wlln}
Let $X_1, X_2, \ldots$ be i.i.d.\ random variables with finite mean
$\mu = \E[X_1]$ and finite variance $\sigma^2 = \Var(X_1)$. Define
$\bar{X}_n = \frac{1}{n}\sum_{i=1}^n X_i$. Then
\[
    \bar{X}_n \cprob \mu,
\]
i.e., for every $\varepsilon > 0$,
$\lim_{n \to \infty} \Prob(|\bar{X}_n - \mu| \geq \varepsilon) = 0$.
\end{theorem}

\begin{proof}
By Example~\ref{ex:chebyshev_mean} (Chebyshev's inequality applied
to $\bar{X}_n$),
\[
    \Prob(|\bar{X}_n - \mu| \geq \varepsilon)
    \leq \frac{\Var(\bar{X}_n)}{\varepsilon^2}
    = \frac{\sigma^2}{n\varepsilon^2}.
\]
For fixed $\varepsilon > 0$, the right side tends to $0$ as
$n \to \infty$. Therefore $\bar{X}_n \cprob \mu$.
\end{proof}

\begin{keyresult}[title={The Law of Large Numbers}]
The weak law of large numbers guarantees that sample averages
converge in probability to the population mean. This is the
theoretical foundation for Monte Carlo estimation: averaging
i.i.d.\ samples from any distribution with finite variance yields
a consistent estimator of the mean. In RL, this justifies computing
$V^\pi(s) \approx \frac{1}{n}\sum_{i=1}^n G^{(i)}$ where the
$G^{(i)}$ are i.i.d.\ returns from state $s$.
\end{keyresult}

\begin{example}[WLLN for a Biased Coin]
\label{ex:wlln_coin}
Toss a coin with $\Prob(\text{heads}) = p$ independently $n$ times.
Let $X_i = \indicator_{\{\text{heads on toss } i\}}$, so each
$X_i \sim \mathrm{Bernoulli}(p)$ with $\mu = p$ and
$\sigma^2 = p(1-p) \leq 1/4$. The sample mean
$\bar{X}_n = (\text{number of heads})/n$ is the empirical frequency
of heads. By the WLLN, $\bar{X}_n \cprob p$: the observed
frequency converges (in probability) to the true probability.
Chebyshev gives the quantitative bound
$\Prob(|\bar{X}_n - p| \geq \varepsilon) \leq 1/(4n\varepsilon^2)$.
\end{example}

\subsection{Strong Law of Large Numbers}

\begin{theorem}[Strong Law of Large Numbers (SLLN)]
\label{thm:slln}
Let $X_1, X_2, \ldots$ be i.i.d.\ random variables with finite mean
$\mu = \E[X_1]$. Then
\[
    \bar{X}_n \cas \mu.
\]
\end{theorem}

The SLLN strengthens the WLLN in two ways: it requires only finite
mean (not finite variance), and it gives almost sure convergence
rather than convergence in probability. The proof requires
measure-theoretic tools -- specifically the Borel-Cantelli lemma and
truncation arguments -- and is deferred to Phase~01
(Lesson~5, probability foundations).

\begin{warning}[title={WLLN vs.\ SLLN: Finite Variance is Not
    Always Needed}]
The WLLN proof above requires finite variance. This is a
limitation of the Chebyshev-based approach, not of the WLLN itself.
A more refined proof using truncation (see \citet[Section 2.2]{durrett2019})
establishes the WLLN under the weaker condition $\E[|X_1|] < \infty$
alone. The SLLN also requires only $\E[|X_1|] < \infty$.
For the purpose of this course, we state the SLLN without proof
and note that in RL applications, bounded rewards automatically
guarantee finite moments of all orders.
\end{warning}

%% ============================================================
%% SECTION 6: THE CENTRAL LIMIT THEOREM
%% ============================================================

\section{The Central Limit Theorem}
\label{sec:clt}

The law of large numbers tells us that $\bar{X}_n \to \mu$; the
central limit theorem describes the \emph{rate} and \emph{shape} of
the fluctuations around $\mu$.

\begin{theorem}[Central Limit Theorem (CLT)]
\label{thm:clt}
Let $X_1, X_2, \ldots$ be i.i.d.\ random variables with finite mean
$\mu = \E[X_1]$ and finite variance $\sigma^2 = \Var(X_1) > 0$. Then
\[
    \frac{\bar{X}_n - \mu}{\sigma / \sqrt{n}} \cdist
    \mathcal{N}(0, 1),
\]
i.e., $\sqrt{n}(\bar{X}_n - \mu)/\sigma$ converges in distribution
to a standard normal random variable.
\end{theorem}

The proof uses characteristic functions (or equivalently, MGFs when
they exist) and is deferred to Phase~01. We focus here on
understanding and applying the CLT.

\begin{remark}
\label{rem:clt_interpretation}
The CLT says that for large $n$, the sample mean is approximately
normal:
\[
    \bar{X}_n \;\dot\sim\; \mathcal{N}\!\left(\mu,
    \frac{\sigma^2}{n}\right).
\]
This approximation improves as $n$ increases and holds regardless of
the distribution of the $X_i$ (provided $\sigma^2$ is finite).
\end{remark}

\begin{example}[CLT for Polling]
\label{ex:clt_polling}
A pollster surveys $n = 1000$ voters, each independently supporting
candidate $A$ with probability $p = 0.52$. Let $\hat{p} = \bar{X}_n$
be the sample proportion. By the CLT,
\[
    \hat{p} \;\dot\sim\; \mathcal{N}\!\left(0.52,
    \frac{0.52 \times 0.48}{1000}\right)
    = \mathcal{N}(0.52, 0.0002496).
\]
The standard deviation of $\hat{p}$ is
$\sigma/\sqrt{n} = \sqrt{0.0002496} \approx 0.0158$.
The probability that the poll shows candidate $A$ trailing
($\hat{p} < 0.5$) is approximately
\[
    \Prob(\hat{p} < 0.5) \approx \Prob\!\left(Z < \frac{0.5 - 0.52}{0.0158}\right)
    = \Prob(Z < -1.27) \approx 0.102,
\]
where $Z \sim \mathcal{N}(0,1)$.
\end{example}

\begin{example}[CLT for Dice Sums]
\label{ex:clt_dice}
Roll a fair six-sided die $n = 100$ times and let
$S_{100} = \sum_{i=1}^{100} X_i$ be the total. Each $X_i$ has
$\mu = 7/2$ and $\sigma^2 = 35/12$ (Lesson~10). By the CLT,
\[
    S_{100} \;\dot\sim\; \mathcal{N}(350, 100 \cdot 35/12)
    = \mathcal{N}(350, 291.67).
\]
The probability that $S_{100}$ exceeds $375$ is approximately
\[
    \Prob(S_{100} > 375) \approx \Prob\!\left(Z > \frac{375 - 350}{\sqrt{291.67}}\right)
    = \Prob(Z > 1.464) \approx 0.072.
\]
Despite the discrete, non-Gaussian distribution of each $X_i$, the
sum is well-approximated by a normal distribution.
\end{example}

\begin{rlconnection}[title={CLT and Confidence Intervals for
    Policy Evaluation}]
In Monte Carlo policy evaluation, the CLT provides confidence
intervals for the estimated value function. If $G^{(1)}, \ldots,
G^{(n)}$ are i.i.d.\ returns with mean $V^\pi(s)$ and variance
$\sigma^2$, then by the CLT an approximate $95\%$ confidence
interval for $V^\pi(s)$ is
\[
    \bar{G}_n \pm 1.96 \cdot \frac{\hat{\sigma}}{\sqrt{n}},
\]
where $\hat{\sigma}$ is the sample standard deviation. This is how
practitioners assess the precision of value estimates and decide
when to stop collecting episodes.
\end{rlconnection}

%% ============================================================
%% SECTION 7: TAIL BOUNDS AND SAMPLE COMPLEXITY
%% ============================================================

\section{Tail Bounds and Sample Complexity}
\label{sec:tail_bounds}

The Chebyshev bound provides a polynomial ($1/n$) rate of
concentration. For bounded random variables, we can obtain much
tighter exponential bounds via the Chernoff method and Hoeffding's
inequality.

\subsection{Chebyshev Confidence Interval}

Before developing Hoeffding's inequality, let us formalize the sample
complexity implied by Chebyshev's inequality.

\begin{proposition}[Chebyshev Sample Complexity]
\label{prop:chebyshev_sample}
Let $X_1, \ldots, X_n$ be i.i.d.\ with mean $\mu$ and variance
$\sigma^2$. To ensure
$\Prob(|\bar{X}_n - \mu| \geq \varepsilon) \leq \delta$,
it suffices to take
\[
    n \geq \frac{\sigma^2}{\varepsilon^2 \delta}.
\]
\end{proposition}

\begin{proof}
From Chebyshev's inequality
(Example~\ref{ex:chebyshev_mean}),
$\Prob(|\bar{X}_n - \mu| \geq \varepsilon) \leq \sigma^2/(n\varepsilon^2)$.
Setting $\sigma^2/(n\varepsilon^2) \leq \delta$ and solving for $n$
gives $n \geq \sigma^2/(\varepsilon^2 \delta)$.
\end{proof}

\begin{example}[Chebyshev Sample Size]
\label{ex:chebyshev_sample}
For i.i.d.\ Bernoulli rewards with $\sigma^2 \leq 1/4$, to estimate
the mean within $\varepsilon = 0.05$ with probability at least
$1 - \delta = 0.95$ ($\delta = 0.05$), Chebyshev requires
\[
    n \geq \frac{1/4}{0.05^2 \cdot 0.05} = \frac{0.25}{0.000125} = 2000.
\]
\end{example}

\subsection{The Chernoff Method}

The key idea for obtaining exponential tail bounds is the
\emph{Chernoff method}: apply Markov's inequality to $e^{tX}$ for
an optimally chosen $t > 0$.

\begin{proposition}[Chernoff Bound -- Upper Tail]
\label{prop:chernoff}
For any random variable $X$ and $a \in \R$,
\[
    \Prob(X \geq a) \leq \inf_{t > 0} e^{-ta} M_X(t),
\]
provided the MGF $M_X(t)$ exists.
\end{proposition}

\begin{proof}
For any $t > 0$, the event $\{X \geq a\}$ is equivalent to
$\{e^{tX} \geq e^{ta}\}$. Since $e^{tX} \geq 0$, Markov's
inequality gives
\[
    \Prob(X \geq a) = \Prob(e^{tX} \geq e^{ta})
    \leq \frac{\E[e^{tX}]}{e^{ta}} = e^{-ta} M_X(t).
\]
This holds for every $t > 0$, so we can take the infimum over $t > 0$.
\end{proof}

\subsection{Hoeffding's Lemma}

\begin{lemma}[Hoeffding's Lemma]
\label{lem:hoeffding}
Let $X$ be a random variable with $\E[X] = 0$ and $a \leq X \leq b$
almost surely. Then for all $t \in \R$,
\[
    \E[e^{tX}] \leq \exp\!\left(\frac{t^2(b-a)^2}{8}\right).
\]
\end{lemma}

\begin{proof}
Since $a \leq X \leq b$, the random variable $X$ can be written as
a convex combination of the endpoints. For any $x \in [a, b]$,
\[
    x = \frac{b - x}{b - a} \cdot a + \frac{x - a}{b - a} \cdot b.
\]
Since $e^{tx}$ is convex in $x$ (it is the exponential function
composed with a linear function), we have
\[
    e^{tx} \leq \frac{b - x}{b - a} e^{ta} + \frac{x - a}{b - a} e^{tb}.
\]
Taking expectations (and using $\E[X] = 0$):
\begin{align}
    \E[e^{tX}] &\leq \frac{b - \E[X]}{b - a} e^{ta}
    + \frac{\E[X] - a}{b - a} e^{tb} \nonumber \\
    &= \frac{b}{b - a} e^{ta} + \frac{-a}{b - a} e^{tb}. \label{eq:hoeffding_step1}
\end{align}
Define $p = -a/(b-a)$ (note $p \in [0,1]$ since $a \leq 0 \leq b$,
which follows from $\E[X] = 0$ and $a \leq X \leq b$) and
$u = t(b-a)$. Then \eqref{eq:hoeffding_step1} becomes
\[
    \E[e^{tX}] \leq (1-p) e^{-pu} + p \, e^{(1-p)u}
    = e^{\varphi(u)},
\]
where $\varphi(u) = -pu + \ln(1 - p + pe^u)$.

We claim that $\varphi(u) \leq u^2/8$ for all $u \in \R$. To
verify this, note $\varphi(0) = 0$ and $\varphi'(0) = 0$:
\[
    \varphi'(u) = -p + \frac{pe^u}{1 - p + pe^u},
    \qquad
    \varphi'(0) = -p + \frac{p}{1} = 0.
\]
The second derivative is
\[
    \varphi''(u) = \frac{pe^u(1 - p + pe^u) - (pe^u)^2}{(1-p+pe^u)^2}
    = \frac{pe^u(1-p)}{(1-p+pe^u)^2}.
\]
Let $q = pe^u / (1 - p + pe^u) \in (0,1)$. Then
$\varphi''(u) = q(1-q) \leq 1/4$
(since $q(1-q)$ achieves its maximum of $1/4$ at $q = 1/2$).

By Taylor's theorem with Lagrange remainder, there exists $\xi$
between $0$ and $u$ such that
\[
    \varphi(u) = \varphi(0) + \varphi'(0)u + \frac{1}{2}\varphi''(\xi)u^2
    = \frac{1}{2}\varphi''(\xi)u^2
    \leq \frac{u^2}{8}.
\]
Since $u = t(b-a)$, we conclude
$\E[e^{tX}] \leq e^{t^2(b-a)^2/8}$.
\end{proof}

\subsection{Hoeffding's Inequality}

\begin{theorem}[Hoeffding's Inequality]
\label{thm:hoeffding}
Let $X_1, \ldots, X_n$ be independent random variables with
$a_i \leq X_i \leq b_i$ almost surely for each $i$. Let
$S_n = \sum_{i=1}^n X_i$ and $\mu = \E[S_n]$. Then for any
$t > 0$,
\[
    \Prob(S_n - \mu \geq t)
    \leq \exp\!\left(-\frac{2t^2}{\sum_{i=1}^n (b_i - a_i)^2}\right)
\]
and
\[
    \Prob(|S_n - \mu| \geq t)
    \leq 2\exp\!\left(-\frac{2t^2}{\sum_{i=1}^n (b_i - a_i)^2}\right).
\]
\end{theorem}

\begin{proof}
Define $Y_i = X_i - \E[X_i]$, so $\E[Y_i] = 0$ and
$a_i - \E[X_i] \leq Y_i \leq b_i - \E[X_i]$. The width of the
interval for $Y_i$ is still $b_i - a_i$. Also,
$S_n - \mu = \sum_{i=1}^n Y_i$.

For any $s > 0$, by the Chernoff method
(Proposition~\ref{prop:chernoff}):
\begin{align*}
    \Prob(S_n - \mu \geq t)
    &= \Prob\!\left(e^{s(S_n - \mu)} \geq e^{st}\right)
    \leq e^{-st} \, \E\!\left[e^{s \sum_{i=1}^n Y_i}\right].
\end{align*}
Since the $Y_i$ are independent,
\[
    \E\!\left[e^{s \sum_i Y_i}\right]
    = \prod_{i=1}^n \E[e^{sY_i}].
\]
By Hoeffding's lemma (Lemma~\ref{lem:hoeffding}), each factor
satisfies $\E[e^{sY_i}] \leq \exp(s^2(b_i - a_i)^2/8)$.
Therefore
\[
    \Prob(S_n - \mu \geq t)
    \leq e^{-st} \prod_{i=1}^n \exp\!\left(\frac{s^2(b_i-a_i)^2}{8}\right)
    = \exp\!\left(-st + \frac{s^2}{8}\sum_{i=1}^n (b_i-a_i)^2\right).
\]
Minimizing over $s > 0$: the exponent
$h(s) = -st + \frac{s^2}{8}\sum_i (b_i - a_i)^2$ is a quadratic in
$s$ with minimum at
\[
    s^* = \frac{4t}{\sum_{i=1}^n (b_i - a_i)^2}.
\]
Substituting:
\[
    h(s^*) = -\frac{4t^2}{\sum_i(b_i-a_i)^2}
    + \frac{16t^2}{8(\sum_i(b_i-a_i)^2)}
    = -\frac{4t^2}{\sum_i(b_i-a_i)^2}
    + \frac{2t^2}{\sum_i(b_i-a_i)^2}
    = -\frac{2t^2}{\sum_i(b_i-a_i)^2}.
\]
This gives the one-sided bound. For the two-sided bound, apply the
same argument to $-Y_i$ (which has the same width) to get
$\Prob(S_n - \mu \leq -t) \leq
\exp(-2t^2/\sum_i(b_i-a_i)^2)$,
then use the union bound.
\end{proof}

\begin{corollary}[Hoeffding for i.i.d.\ Bounded Variables]
\label{cor:hoeffding_iid}
If $X_1, \ldots, X_n$ are i.i.d.\ with $a \leq X_i \leq b$ and
$\mu = \E[X_1]$, then for the sample mean
$\bar{X}_n = \frac{1}{n}\sum_{i=1}^n X_i$,
\[
    \Prob(|\bar{X}_n - \mu| \geq \varepsilon)
    \leq 2\exp\!\left(-\frac{2n\varepsilon^2}{(b-a)^2}\right).
\]
\end{corollary}

\begin{proof}
Apply Theorem~\ref{thm:hoeffding} to $S_n = n\bar{X}_n$ with
$t = n\varepsilon$ and $b_i - a_i = b - a$ for all $i$:
\[
    \Prob(|n\bar{X}_n - n\mu| \geq n\varepsilon)
    \leq 2\exp\!\left(-\frac{2n^2\varepsilon^2}{n(b-a)^2}\right)
    = 2\exp\!\left(-\frac{2n\varepsilon^2}{(b-a)^2}\right). \qedhere
\]
\end{proof}

\subsection{Sample Complexity via Hoeffding}

\begin{proposition}[Hoeffding Sample Complexity]
\label{prop:hoeffding_sample}
Let $X_1, \ldots, X_n$ be i.i.d.\ with $a \leq X_i \leq b$. To
ensure $\Prob(|\bar{X}_n - \mu| \geq \varepsilon) \leq \delta$,
it suffices to take
\[
    n \geq \frac{(b-a)^2}{2\varepsilon^2} \ln\frac{2}{\delta}.
\]
\end{proposition}

\begin{proof}
By Corollary~\ref{cor:hoeffding_iid}, we need
$2\exp(-2n\varepsilon^2/(b-a)^2) \leq \delta$. Solving:
\[
    \exp\!\left(-\frac{2n\varepsilon^2}{(b-a)^2}\right)
    \leq \frac{\delta}{2},
    \quad\text{so}\quad
    \frac{2n\varepsilon^2}{(b-a)^2} \geq \ln\frac{2}{\delta},
    \quad\text{i.e.,}\quad
    n \geq \frac{(b-a)^2}{2\varepsilon^2}\ln\frac{2}{\delta}. \qedhere
\]
\end{proof}

\begin{example}[Hoeffding vs.\ Chebyshev Sample Size]
\label{ex:hoeffding_vs_chebyshev}
Consider the same setting as Example~\ref{ex:chebyshev_sample}:
i.i.d.\ Bernoulli rewards with $a = 0$, $b = 1$,
$\varepsilon = 0.05$, and $\delta = 0.05$. Hoeffding requires
\[
    n \geq \frac{(1-0)^2}{2 \cdot 0.05^2}\ln\frac{2}{0.05}
    = \frac{1}{0.005} \ln(40)
    = 200 \cdot 3.689 \approx 738.
\]
Compare with Chebyshev's requirement of $n \geq 2000$. Hoeffding
gives a nearly $2.7\times$ reduction in sample size. This
improvement comes from exploiting the bounded range of the variable,
not just its variance.
\end{example}

\begin{keyresult}[title={Sample Complexity: Chebyshev vs.\ Hoeffding}]
For estimating the mean of bounded i.i.d.\ random variables
$X_i \in [a, b]$ within accuracy $\varepsilon$ and confidence
$1 - \delta$:
\begin{center}
\begin{tabular}{lll}
\toprule
\textbf{Method} & \textbf{Sample complexity} & \textbf{Dependence on $\delta$} \\
\midrule
Chebyshev & $n \geq \sigma^2 / (\varepsilon^2 \delta)$
    & Polynomial: $O(1/\delta)$ \\
Hoeffding & $n \geq \frac{(b-a)^2}{2\varepsilon^2}\ln(2/\delta)$
    & Logarithmic: $O(\ln(1/\delta))$ \\
\bottomrule
\end{tabular}
\end{center}
For high-confidence estimates (small $\delta$), Hoeffding's
logarithmic dependence on $1/\delta$ vastly outperforms Chebyshev's
polynomial dependence. This is why exponential tail bounds are
essential in PAC learning theory.
\end{keyresult}

\begin{rlconnection}[title={Monte Carlo Policy Evaluation Sample
    Complexity}]
Consider estimating $V^\pi(s)$ by averaging $n$ i.i.d.\ returns
$G^{(1)}, \ldots, G^{(n)}$. If rewards are bounded in $[0, R_{\max}]$
and the discount factor is $\gamma < 1$, each return satisfies
$0 \leq G^{(i)} \leq R_{\max}/(1 - \gamma)$. By
Proposition~\ref{prop:hoeffding_sample},
\[
    n \geq \frac{R_{\max}^2}{2(1-\gamma)^2 \varepsilon^2}
    \ln\frac{2}{\delta}
\]
episodes suffice to ensure
$|\bar{G}_n - V^\pi(s)| < \varepsilon$ with probability at least
$1 - \delta$. This is the basic sample complexity of Monte Carlo
policy evaluation. Note the dependence on $(1-\gamma)^{-2}$: as the
discount factor approaches $1$, the effective horizon grows and far
more episodes are needed.
\end{rlconnection}

%% ============================================================
%% SECTION 8: EXERCISES
%% ============================================================

\section{Exercises}
\label{sec:exercises}

\begin{exercise}[MGF Computations]
\textnormal{(\(\ast\))} \\
Let $X \sim \mathrm{Geometric}(p)$ with
$\Prob(X = k) = (1-p)^{k-1}p$ for $k = 1, 2, \ldots$
\begin{enumerate}[nosep,label=(\alph*)]
    \item Compute the MGF $M_X(t)$ and determine the values of $t$
    for which it is finite.
    \item Use the MGF to compute $\E[X]$ and $\Var(X)$.
\end{enumerate}
\end{exercise}

\begin{exercise}[Jensen's Inequality Applications]
\textnormal{(\(\ast\))} \\
Let $X$ be a positive random variable with finite expectation.
\begin{enumerate}[nosep,label=(\alph*)]
    \item Use Jensen's inequality with $g(x) = 1/x$ (which is convex
    on $(0, \infty)$) to prove that $\E[1/X] \geq 1/\E[X]$.
    \item Use Jensen with $g(x) = -\ln x$ to prove that
    $\E[\ln X] \leq \ln(\E[X])$.
    \item Let $X \sim \mathrm{Uniform}(1, 3)$. Verify part (a)
    numerically by computing both $\E[1/X]$ and $1/\E[X]$.
\end{enumerate}
\end{exercise}

\begin{exercise}[Chebyshev Sharpness]
\textnormal{(\(\ast\ast\))} \\
Construct a random variable $X$ with $\E[X] = 0$ and $\Var(X) = 1$
such that $\Prob(|X| \geq 2) = 1/4$, showing that Chebyshev's
bound $\Prob(|X| \geq 2) \leq 1/4$ is tight for $k = 2$.
\\
\textit{Hint:} Consider a three-point distribution.
\end{exercise}

\begin{exercise}[WLLN Does Not Require Identical Distributions]
\textnormal{(\(\ast\ast\))} \\
Let $X_1, X_2, \ldots$ be independent (but not necessarily
identically distributed) random variables with $\E[X_i] = \mu$ for
all $i$ and $\Var(X_i) \leq C$ for some constant $C < \infty$.
Use Chebyshev's inequality to prove that
$\bar{X}_n = \frac{1}{n}\sum_{i=1}^n X_i \cprob \mu$.
\end{exercise}

\begin{exercise}[Monte Carlo Sample Complexity]
\textnormal{(\(\ast\ast\))} \\
An RL agent estimates $V^\pi(s)$ using Monte Carlo rollouts. Rewards
are bounded in $[0, 1]$ and $\gamma = 0.99$.
\begin{enumerate}[nosep,label=(\alph*)]
    \item Compute the range $[0, R_{\max}/(1-\gamma)]$ of a single
    discounted return.
    \item Using Hoeffding's inequality, how many episodes are needed
    to ensure the estimate is within $\varepsilon = 0.5$ of the true
    value with probability at least $0.95$?
    \item How does the answer change if $\gamma = 0.9$?
    \item Compare the Hoeffding sample complexity with the Chebyshev
    sample complexity, using the bound
    $\Var(G) \leq (b-a)^2/4$.
\end{enumerate}
\end{exercise}

\begin{exercise}[Proving Hoeffding via Symmetrization]
\textnormal{(\(\ast\ast\ast\))} \\
Let $X$ be a random variable with $\E[X] = 0$ and $a \leq X \leq b$.
An alternative route to Hoeffding's lemma uses the following steps:
\begin{enumerate}[nosep,label=(\alph*)]
    \item Let $X'$ be an independent copy of $X$. Show that
    $\E[e^{t(X - X')}] \leq \E[e^{tX}]^2$.

    \textit{Hint:} Use Jensen's inequality applied to the conditional
    expectation $\E[e^{-tX'} \mid X]$.
    \item The symmetric variable $X - X'$ satisfies
    $|X - X'| \leq b - a$. Using the fact that a symmetric bounded
    random variable $Y$ with $|Y| \leq c$ satisfies
    $\E[e^{tY}] \leq e^{t^2 c^2 / 2}$ (which can be shown via the
    Taylor series of $\cosh$), conclude that
    $\E[e^{tX}] \leq e^{t^2(b-a)^2/8}$.
\end{enumerate}
\end{exercise}

%% ============================================================
%% SECTION 9: SUMMARY
%% ============================================================

\section{Summary}
\label{sec:summary}

\begin{itemize}
    \item The \emph{moment generating function} $M_X(t) = \E[e^{tX}]$
    encodes all moments via $\E[X^n] = M_X^{(n)}(0)$ and, when it
    exists in a neighborhood of the origin, uniquely determines the
    distribution.
    \item The MGF of a sum of independent random variables factorizes:
    $M_{X+Y}(t) = M_X(t) M_Y(t)$. This immediately shows that sums
    of independent Poissons are Poisson and sums of independent
    normals are normal.
    \item \emph{Markov's inequality}
    $\Prob(X \geq a) \leq \E[X]/a$ requires only non-negativity and
    a finite first moment.
    \item \emph{Chebyshev's inequality}
    $\Prob(|X - \mu| \geq k) \leq \sigma^2/k^2$ uses the first
    two moments and gives polynomial tail decay.
    \item \emph{Cauchy-Schwarz}
    $|\E[XY]|^2 \leq \E[X^2]\E[Y^2]$ and \emph{Jensen's inequality}
    $g(\E[X]) \leq \E[g(X)]$ (for convex $g$) are workhorse tools
    throughout probability and optimization.
    \item The \emph{weak law of large numbers} states
    $\bar{X}_n \cprob \mu$ (proved via Chebyshev); the \emph{strong
    law} strengthens this to $\bar{X}_n \cas \mu$ (proof deferred to
    Phase~01).
    \item The \emph{central limit theorem} states
    $\sqrt{n}(\bar{X}_n - \mu)/\sigma \cdist \mathcal{N}(0,1)$:
    sample means are approximately normal for large $n$, regardless
    of the underlying distribution.
    \item \emph{Hoeffding's inequality} provides exponential
    concentration for bounded random variables:
    $\Prob(|\bar{X}_n - \mu| \geq \varepsilon) \leq
    2e^{-2n\varepsilon^2/(b-a)^2}$.
    \item The Hoeffding sample complexity
    $n \geq \frac{(b-a)^2}{2\varepsilon^2}\ln(2/\delta)$
    depends logarithmically on $1/\delta$, vastly outperforming
    Chebyshev's polynomial dependence for high-confidence estimation.
\end{itemize}

\bigskip

As the capstone lesson of the Probability subdomain, we now review
the arc of Lessons~9--12.

\begin{center}
\begin{tabular}{@{}clp{7.5cm}@{}}
\toprule
\textbf{Lesson} & \textbf{Title} & \textbf{Core Contribution} \\
\midrule
9 & Probability Spaces and Conditioning
  & Kolmogorov axioms, conditional probability, Bayes' theorem,
    independence. Provides the foundational language. \\
10 & Random Variables and Distributions
  & Random variables, PMF/PDF/CDF, Bernoulli, Binomial, Geometric,
    Poisson, Uniform, Exponential, Normal distributions.
    Introduces the objects we compute with. \\
11 & Expectation, Variance, and Joint Distributions
  & Linearity of expectation, LOTUS, variance, covariance, joint
    and conditional distributions, tower property, law of total
    variance. Provides the computational machinery. \\
12 & Inequalities and Limit Theorems
  & MGFs, Markov/Chebyshev/Cauchy-Schwarz/Jensen inequalities,
    WLLN, CLT, Hoeffding's inequality, sample complexity bounds.
    Provides the asymptotic and finite-sample guarantees. \\
\bottomrule
\end{tabular}
\end{center}

%% ============================================================
%% SECTION 10: FURTHER READING
%% ============================================================

\section{Further Reading}
\label{sec:further_reading}

\begin{itemize}
    \item \textbf{Phase 01, Lesson 4} (Measure Spaces) provides the
    measure-theoretic foundations that make the definitions in
    Lessons~9--12 fully rigorous, replacing finite sums and integrals
    with Lebesgue integration.
    \item \textbf{Phase 01, Lesson 5} (Probability Foundations)
    revisits convergence concepts, proves the strong law of large
    numbers using the Borel-Cantelli lemma, establishes the CLT via
    characteristic functions, and develops martingale theory.
    \item \textbf{Phase 01, Lesson 5} also covers the
    \emph{Azuma-Hoeffding inequality} for martingale difference
    sequences, which generalizes Hoeffding's inequality to dependent
    random variables -- essential for analyzing online RL algorithms
    where successive observations are correlated.
    \item \textbf{Phase 02} (Core RL) uses the law of large numbers
    to justify Monte Carlo estimation, Chebyshev and Hoeffding
    inequalities for finite-sample convergence bounds, and the CLT
    for asymptotic confidence intervals in policy evaluation.
    \item For a comprehensive treatment of concentration inequalities
    beyond Hoeffding (including Bernstein, Bennett, and
    McDiarmid's inequalities), see
    \citet[Chapter 2]{durrett2019} and \citet[Appendix A]{grimmett2020}.
    \item The Chernoff bound technique (optimizing over $t$ in
    Markov's inequality applied to $e^{tX}$) is a special case of
    the \emph{exponential tilting} method, which plays a central
    role in large deviations theory; see \citet[Chapter 8]{feller1968}
    for an introduction.
    \item Historically, the law of large numbers was first proved by
    Jakob Bernoulli (1713) in a special case (Bernoulli trials).
    Chebyshev (1867) gave the first general proof using
    his inequality. The central limit theorem was established by
    Lindeberg (1922) and Feller (1935) under progressively weaker
    conditions.
\end{itemize}

\bibliography{probability-refs}

\end{document}
